\documentclass[12pt,a4paper,oneside]{report}
% nagyon sok kép esetén meggyorsítható a fordítás a draft móddal
% \documentclass[12pt,a4paper,oneside,draft]{report}
% ekkor a képek nem renderelődnek ki, csak placeholder lesz mérethelyesen
\usepackage[utf8]{inputenc} % mindenképp maradjon az utf-8 kódolás
\usepackage[magyar]{babel}
\usepackage[T1]{fontenc}
\usepackage{amsmath}
\usepackage{amsfonts}
\usepackage{amssymb}
\usepackage{graphics} % grafikus elemek, képek berakásához
\usepackage{epsfig} % eps importáláshoz
\usepackage{listings}
\usepackage{sectsty}
\usepackage{enumerate}
\usepackage{lastpage}
\usepackage{setspace}
\usepackage{hyperref} % PDF hivatkozásokhoz kell
\usepackage[hang]{caption}
\usepackage{titling} % a title, author parancsok szabad használatához

% a TikZ rajzoló modul, és a kapcsolási rajz készítő modul, ha kell
%\usepackage{tikz}
%\usepackage{circuitikz}

% az A4 oldal margóinak és méreteinek beállítása
\usepackage[left=25mm,right=25mm,top=20mm,bottom=25mm]{geometry}\pagestyle{plain}

% A sorköz távolság beállítása
% egyszeres sorköz
\singlespacing
% 1,5 sorköz
% \onehalfspacing

% A hivatkozások, és linkek átállítása alapértelmezett színre, fekete-fehér nyomtatáshoz optimalizálva
\hypersetup
{
  	colorlinks,
  	citecolor=black,
 	linkcolor=black,
  	urlcolor=black
}

% --------- A hallgató által kitöltendő rész! --------- %

% a dolgozat típusa
\title{Önálló laboratórium beszámoló} 

% a dolgozat szerzője
\author{Bálint Éliás}

% a dolgozat típusát adjuk meg tárgyesetben, kisbetűvel (pl.: diplomatervet, szakdolgozatot, témalabor beszámolót stb.)
\newcommand{\dokumentumtipus}{önálló labor beszámolót} 

% a dolgozat témája
\newcommand{\dokumentumcim}{Szinkronizációs eljárások vizsgálata 5G OFDMA rendszerekben IEEE specifikációk alapján}

% --------- A hallgató által kitöltendő rész! --------- %

\begin{document}
\begin{titlepage}

\begin{figure}
\centering
\includegraphics[width=100mm,keepaspectratio]{bme.pdf}
\end{figure}

\centering
\textbf{Budapesti Műszaki és Gazdaságtudományi Egyetem}\\
\textbf{Villamosmérnöki és Informatikai Kar}\\
\textbf{Szélessávú Hírközlés és Villamosságtan Tanszék}\\
\vspace{5mm}
\includegraphics[width=40mm,keepaspectratio]{hvt_logo_only_fixed_vector_inverted.png}  \\
\vspace{40mm}
\Huge
\dokumentumcim\\
\vspace{50mm}
\Large
\thetitle \\
\vspace{12mm}
\Large
\textbf{\theauthor}\\
\vspace{12mm}
\the\year


\end{titlepage}
% \newpage
\hspace{10mm}
\begin{center}
\large
\textbf{HALLGATÓI NYILATKOZAT}
\end{center}
\hspace{10mm}

Alulírott \theauthor, szigorló hallgató kijelentem, hogy ezt a \dokumentumtipus \ meg nem engedett segítség nélkül, saját magam készítettem, csak a megadott forrásokat (szakirodalom, eszközök stb.) használtam fel. Minden olyan részt, melyet szó szerint, vagy azonos értelemben, de átfogalmazva más forrásból átvettem, egyértelműen, a forrás megadásával megjelöltem.

Hozzájárulok, hogy a jelen munkám alapadatait (szerző(k), cím, angol és magyar nyelvű tartalmi kivonat, készítés éve, konzulens(ek) neve) a BME VIK nyilvánosan hozzáférhető elektronikus formában, a munka teljes szövegét pedig az egyetem belső hálózatán keresztül (vagy autentikált felhasználók számára) közzétegye. Kijelentem, hogy a benyújtott munka és annak elektronikus verziója megegyezik. Dékáni engedéllyel titkosított diplomatervek esetén a dolgozat szövege csak 3 év eltelte után válik hozzáférhetővé.

\hspace{10mm}

Budapest, \today

\begin{flushright}
\theauthor
\end{flushright}
% Ide kell include-olni, vagy befűzni a feladatkiírást, ha van
\tableofcontents
% \chapter*{Kivonat}

Jelen dokumentum egy diplomaterv sablon, amely formai keretet ad a BME Villamosmérnöki és Informatikai Karán végző hallgatók által elkészítendő szakdolgozatnak, és diplomatervnek. A sablon használata opcionális. Ez a sablon \LaTeX \  alapú, a TeXLive TeX 
implementációval és a PDF-LaTeX fordítóval működőképes.

% \chapter*{Abstract}

This document is a \LaTeX -based template for the BSc/MSc thesis of students at the Electrical Engineering and Informatics Faculty of Budapest University of Technology and Economics. The usage of this template is optional. It has been tested with the TeXLive TEX implementation, and it requires the PDF-LaTeX compiler.
\chapter{Bevezetés}

\section{Előző féléves munkám}

A Budapesti Műszaki és Gazdaságtudományi Egyetem Villamosmérnöki és Informatikai karán a villamosmérnök alapszak 5. félévében Témalaboratórium tantárgyat hirdetnek.
Ez a tantárgy egy olyan projekttárgy, amely a hallgatók számára lehetővé teszi, hogy valamilyen területen elmélyedjenek, és közösen dolgozzanak egy témavezetővel, ezzel tapasztalatot nyújtva számukra, a későbbi Önálló laboratórium és Szakdolgozat készítés tárgyakhoz.

Én az 5. félévemben a Témalaboratórium keretei között Dr. Horváth Bálint Péter témavezetésével folytattam egyéni munkát, amely során egy OFDM (Orthogonal Frequency Division Multiplexing - Ortogonális Frekvenciaosztásos Nyalábolás) rendszert modelleztem. A munka során Python nyelvben készítettem el egy digitális jelátvitel rendszer modelljét, amelyben az adót, a vevőt, és köztük az átviteli csatorna volt a 3 legfőbb elem. Munkám során azt vizsgáltam, hogy miképp tud egy ilyen OFDM rendszer védekezni a többutas terjedés problémái ellen, milyen lehetőségek vannak a vett jel kompenzálására, vagy az átvitel minőségének mérésére.

Az eredmények részletes tárgyalása nem része ennek a dolgozatnak, de az eredményeket is és a tanultakat is felhasználtam a féléves munkámban.

\section{OFDM és OFDMA}

Az OFDM és az OFDMA (Orthogonal Frequency Division Multiple Access) rendszerek között egy apró különbség van. Míg egy OFDM rendszerben egy adó és egy vevő közötti kommunikáció valósul meg több vivőfrekvencián, addig az OFDMA rendszerben a kommunikáció egy vagy több adó és egy vagy több vevő között valósul meg. Ez számomra nem volt evidens, a tanulmányaim során eddig a kettőt összemosták.

Az OFDM előnyeit könnyedén kihasználhatjuk többszörös hozzáférési rendszerekben is.
Így nevezzük azokat a rendszereket, amelyekben egyszerre párhuzamosan több jelátvitel valósul meg.
Ezt megtehetjük a frekvenciatartománybeli szétválasztással, lásd a GSM uplink és downlink kapcsolat, ahol két teljesen elkülönülő frekvenciatartományt vesz igénybe a kétfajta átvitel. 
Ugyanez a hatás elérhető, ha az átviteli médiumot valamilyen időosztásban osztjuk fel a felhasználók között. 
Van olyan módszer is, amellyel időben és frekvenciában is egyszerre oldhatjuk meg az ütközéseket, ha megfelelően kódoljuk az átvinni kívánt jelet.

\section{OFDMA kihívásai}

xdxd

\section{Munkám}
% \chapter*{Köszönetnyilvánítás}
Ide kerül az opcionális köszönetnyilvánítás, amelyben érdemes megemlíteni a konzulenst, illetve mindazokat, akik a dolgozat létrejöttében segítettek a szerzőnek. 
\begin{thebibliography}{4}

\bibitem{hvthonlap} \url{http://hvt.bme.hu}

\bibitem{radiotechnika} Rádiótechnika évkönyve 2007, 172. oldaltól

\bibitem{diploma} Dudás Levente, \emph{Digitális nyalábformálású antenna (DBF)} diplomaterv, BME SzHVT, 2007

\bibitem{kicad} \url{http://kicad-pcb.org/}

\end{thebibliography}


% \listoffigures
% \listoftables
\chapter*{Függelék}

A függelék az a fejezet, amely nem képezi szerves részét magának a dolgozatnak (nem növeli az oldalszámot), csak a megértést segíti, illetve a plusz ábrákat, pl. kapcsolási rajzokat, NYÁK terveket, program forráskódokat tartalmazza, de természetesen lehet hivatkozni rá (és sok esetben kell is). \label{fugg}

Jelen dokumentumot, nemcsak pdf, hanem \LaTeX \  forráskód szinten is közzéteszem, abból a megfontolásból, hogy lehetőség szerint legyen egységes a beadott beszámolók és jegyzőkönyvek formátuma és kinézete.

Jelen dokumentum forráskódja példákat tartalmaz a szövegek, képek, táblázatok, számozott és számozatlan felsorolások szerkesztésére, tagolására, formátumára vonatkozóan beleértve a matematika kifejezések forráskód szintű kezelését is.

A LaTeX letölthető többféle operációs rendszerre innen: 
\\
\url{https://www.latex-project.org/get/}

Ha valakinek magyar nyelvű szótárra van szüksége, akkor használhatja jelen dokumentum mappájában levő állományt is: a fordítóban kell beállítani a nyelvi beállításoknál.

% végső fordításnál kapcsold ki a draft módot !!

\end{document}